% Source: https://artofproblemsolving.com/community/c3h2579347_question_about_increasing_functions
% Stable identity: aops-topic-2579347-q01
\textbf{(a)} Let $x\ne y$, say $x<y$. Then $f(x)<f(y)$, so $f(x)\ne f(y)$. Hence $f$ is injective. $\blacksquare$

\textbf{(b)} \textbf{No.} The function $f(x)=2^x$ is strictly increasing, but it never takes the value $0$ (or any negative value), so it is not surjective onto $\mathbb R$ and not a bijection. A strictly increasing function is a bijection $\mathbb R\to\mathbb R$ exactly when it is also onto, as for $f(x)=x^3$.

\textbf{(c)} Then even injectivity can fail: the constant function $f(x)=0$ satisfies $f(x)\le f(y)$ for all $x<y$, but it is not injective.
